\chapter{实施过程与效果分析}

\section{实施过程}

第一轮行动研究于2025年10月至11月进行，持续四周。研究者按照设计方案组织学生完成入项、探究、创作与展示四个阶段。第一轮实施总体顺利，学生对走访与访谈表现出较高热情，部分小组能够围绕遗产点提出有价值的问题。例如，有小组在走访洛阳桥后，关注“桥墩为什么做成船形”“古代没有起重机，巨大石板如何架设”等问题，并主动查阅资料、访谈讲解员。也有小组在开元寺观察到东西塔的浮雕与楹联，尝试把佛教故事、建筑细节与泉州多元文化联系起来。

第一轮也暴露出一些问题。其一，任务范围过宽。部分小组选择“泉州世界遗产”作为整体主题，导致资料搜集面过广，解说词缺乏焦点。其二，资料筛选能力不足。学生容易直接复制网络资料，对信息的来源、准确性与适用性缺乏判断。其三，小组合作不均衡。少数学生承担了大部分写作与制作任务，其他成员参与度不高。其四，评价标准理解不到位。学生对“受众意识”“语言得体”等要求感到抽象，修改时不知道从何处入手。

针对上述问题，研究者在第二轮实施前进行了四方面调整。第一，缩小主题范围，要求每组聚焦一处遗产点，并从“历史故事”“人物与生活”“建筑与审美”“保护与传承”四个视角中选择一个切入。第二，增加资料筛选微型课，教学生使用“来源—证据—用途”三栏表格评估资料，并为每组提供两份经过筛选的地方文献选段。第三，引入角色卡与小组公约，明确组长、资料员、访谈员、写作员、展示员等角色，要求角色在项目过程中轮换。第四，将评价量表转化为学生可理解的问题清单，例如“你的解说词是否能让没有去过泉州的人听懂”“你是否用到了访谈或观察中的具体细节”，并在每次修改前进行对照。

第二轮行动研究于2025年11月至12月进行，仍为四周。调整后，小组主题更加集中，资料卡片的质量明显提高，课堂讨论也更有针对性。学生开始主动追问“这个说法有没有依据”“这段材料适合放在开头还是结尾”“我们的解说词是写给谁看的”。在最终展示中，多数小组能够把遗产点信息、现场观察、人物访谈与个人感受结合起来，作品的可读性与现场感较第一轮有明显提升。

\section{数据来源与分析}

本研究的数据包括问卷、课堂观察、学生作品与访谈四类。问卷在前测与后测各发放四十二份，回收有效问卷四十二份，有效率为百分之百。问卷采用五点计分，包含文化理解、语言运用、合作探究与问题解决四个维度，共十六个题项。研究者使用克隆巴赫系数检验问卷内部一致性，结果如式~\eqref{eq:alpha}所示。

\begin{equation}
  \alpha=\frac{k}{k-1}\left(1-\frac{\sum_{i=1}^{k}s_i^2}{s_T^2}\right)
  \label{eq:alpha}
\end{equation}

其中，$k$为题项数，$s_i^2$为第$i$个题项的方差，$s_T^2$为总分的方差。本研究中$k=16$，计算得到$\alpha=0.87$，表明问卷具有较好的内部一致性。课堂观察记录共十六课时，采用叙事记录与关键事件记录相结合的方式。学生作品包括资料卡片、访谈记录、解说词初稿与修改稿、最终成果与反思日志。访谈对象包括十二名学生和四名教师，每次访谈约二十至三十分钟，访谈提纲见附录~\ref{app:interview}。访谈资料经过转写、编码与主题归纳，形成“文化理解”“语言运用”“合作探究”“困难与建议”四个主题。

\section{实施效果}

\subsection{问卷结果}

前后测各维度得分如表~\ref{tab:results}所示。配对样本$t$检验结果显示，四个维度的后测得分均显著高于前测，且效应量达到中等以上水平。其中，文化理解维度的提升最为明显，均值由3.12提高到4.21，效应量为2.02；语言运用维度由3.05提高到3.92，效应量为1.49；合作探究维度由3.48提高到4.16，效应量为1.24；问题解决维度由3.21提高到3.88，效应量为1.10。各维度前后测均值对比如图~\ref{fig:means}所示。

\begin{table}[htbp]
  \centering
  \caption{前后测各维度得分比较（$N=42$）}
  \label{tab:results}
  \begin{tabular}{@{}lcccccc@{}}
    \toprule
    维度 & 前测$M$ & 前测$SD$ & 后测$M$ & 后测$SD$ & $t$ & $d$ \\
    \midrule
    文化理解 & 3.12 & 0.58 & 4.21 & 0.49 & $-9.84^{***}$ & 2.02 \\
    语言运用 & 3.05 & 0.61 & 3.92 & 0.55 & $-8.12^{***}$ & 1.49 \\
    合作探究 & 3.48 & 0.57 & 4.16 & 0.52 & $-7.36^{***}$ & 1.24 \\
    问题解决 & 3.21 & 0.63 & 3.88 & 0.58 & $-6.71^{***}$ & 1.10 \\
    \bottomrule
  \end{tabular}
  \par\smallskip
  {\zihao{6}注：$^{***}p<0.001$。}
\end{table}

\begin{figure}[htbp]
  \centering
  \begin{tikzpicture}
    \begin{axis}[
      ybar,
      bar width=11pt,
      width=0.88\textwidth,
      height=6.2cm,
      ymin=0,
      ymax=5,
      ylabel={平均分},
      symbolic x coords={文化理解,语言运用,合作探究,问题解决},
      xtick=data,
      xticklabel style={font=\zihao{5}},
      yticklabel style={font=\zihao{5}},
      ylabel style={font=\zihao{5}},
      legend style={at={(0.5,1.03)}, anchor=south, legend columns=2,
        font=\zihao{6}, draw=none},
      enlarge x limits=0.16,
      axis lines*=left,
      ymajorgrids=true,
      grid style={dashed, gray!30}
    ]
      \addplot[fill=gray!55, draw=black] coordinates {
        (文化理解,3.12) (语言运用,3.05) (合作探究,3.48) (问题解决,3.21)};
      \addplot[fill=gray!15, draw=black] coordinates {
        (文化理解,4.21) (语言运用,3.92) (合作探究,4.16) (问题解决,3.88)};
      \legend{前测,后测}
    \end{axis}
  \end{tikzpicture}
  \caption{前后测各维度平均分对比}
  \label{fig:means}
\end{figure}

\subsection{作品分析}

研究者依据评价量表对四十二套最终作品进行评分，并从中选取十二套作品进行文本分析。结果显示，多数作品能够做到信息准确、结构完整、语言通顺，部分作品表现出较强的受众意识与创意表达。例如，有小组以“给第一次来泉州的同龄人写一封邀请信”为开头，把洛阳桥的历史、建筑特点与自己的走访体验串联起来；有小组为开元寺设计了一条“从山门到东西塔”的导览路线，在解说词中穿插了弘一法师与泉州佛教文化的故事；还有小组制作了音频导览，用普通话与闽南语双语介绍蟳埔女习俗，尝试保留地方语言的声音质感。

作品分析也发现两类典型问题。第一类是“资料堆砌”，学生把搜集到的信息逐条罗列，缺少取舍与组织，导致解说词像说明书的拼接。第二类是“抒情过度”，学生大量使用“震撼”“伟大”“自豪”等词语，却缺少具体细节与证据支撑，使表达显得空泛。针对这些问题，教师在第二轮中增加了“用细节说话”的写作微课，引导学生从观察记录与访谈中选取最有表现力的细节，并学习用事实支撑感受。

\subsection{访谈结果}

对十二名学生的访谈资料进行编码后，形成三个主要主题。第一个主题是“从陌生到亲近”。多名学生提到，以前觉得世界遗产“离自己很远”，走访之后才发现“原来就在每天经过的地方”。有学生说：“我以前只知道开元寺可以拜拜，不知道东西塔上面有那么多故事。现在带朋友去，我可以讲出很多东西。”第二个主题是“语言表达的变化”。学生普遍认为，写解说词比写普通作文更“难”，因为要考虑“别人能不能听懂”“有没有吸引力”，但也更有成就感。有学生表示：“以前写作文是给老师看的，这次是给来参观的人看的，所以会想怎么开头、怎么讲清楚。”第三个主题是“合作与困难”。学生提到，小组分工、意见分歧和时间安排是项目实施中的主要困难，但通过角色卡、小组会议和教师协调，多数小组能够完成任务。

教师访谈则提供了另一个视角。教师A认为，项目化学习让学生“动起来了”，但“语文课的时间被压缩了”，需要在活动与语言训练之间取得平衡。教师B指出，学生在资料筛选与文本修改方面仍需要更多指导，“不能指望一次项目就能解决所有问题”。教师C建议，学校可以建立地方文化资源库，减少教师搜集资料的时间。教师D则强调，评价不能只看最终作品，否则容易回到“比谁做得漂亮”的老路。这些意见为方案的后续改进提供了重要参考。

\section{讨论}

\subsection{项目化学习何以促进文化理解}

从问卷与访谈结果看，项目化学习对文化理解的促进作用较为明显。这可能与三个因素有关。第一，真实情境激发了学生的探究动机。学生不再是被动接受“泉州是世界遗产城市”这一结论，而是通过走访、观察、访谈与写作，主动建构对遗产的理解。第二，多模态材料提供了丰富的感知通道。图片、视频、建筑实物、地方文献、口述故事与闽南语声音共同构成了学习资源，使学生在看、听、说、写中形成多层面的文化体验。第三，公开成果赋予了学习以社会责任。学生知道自己的作品会被真实受众阅读或观看，因此更愿意投入精力，也更在意表达的准确性与感染力。

不过，文化理解的提升并不自动带来文化认同。学生在访谈中表达的自豪感，有时建立在“泉州很厉害”的笼统印象上，缺少对遗产保护复杂性、社区生活变迁与旅游开发影响的理解。这提示我们，地方遗产教育不能只讲“辉煌”，也要引导学生关注遗产的当代处境，讨论保护与利用、传承与创新之间的关系。只有这样，文化理解才能走向更深层的文化自觉。

\subsection{项目化学习何以促进语言运用}

语言运用维度的提升，与项目中反复出现的真实表达任务密切相关。学生需要根据受众与目的撰写解说词，需要在访谈中提问与追问，需要在小组讨论中陈述理由、回应异议，需要在展示中控制语速、语气与肢体语言。这些任务把抽象的“语言能力”转化为具体可感的表达实践，使学生能够在“用”中学、在“改”中提升。

但研究也发现，语言运用能力的提升并不均衡。部分学生能够写出流畅的解说词，却在口头展示时紧张、念稿；部分学生擅长口头表达，但书面文本结构松散、用词随意。这说明项目化学习需要为不同能力水平的学生提供差异化支架。例如，为写作困难的学生提供句式模板与段落结构图，为口头表达困难的学生提供模拟演练与同伴反馈，为能力较强的学生提供更具挑战性的修改任务。只有在“共同任务”中嵌入“分层支持”，项目才能真正面向全体学生。

\subsection{项目化学习实施中的张力}

项目实施过程中始终存在三组张力。第一是“项目逻辑”与“课时逻辑”的张力。项目化学习需要较长时间持续探究，而学校课程安排以课时为单位，二者之间容易产生冲突。本研究的处理方式是将项目压缩在十二课时内，并将部分探究任务延伸到课外，但这仍然依赖学校的支持与家长的配合。第二是“学生自主”与“教师指导”的张力。项目化学习强调学生的自主选择与探究，但初中生在资料筛选、问题聚焦、文本修改等方面仍需要教师指导。教师既不能包办代替，也不能放任自流，而应在关键节点提供恰到好处的支架。第三是“过程评价”与“结果评价”的张力。过程性证据能够反映学生的学习经历，但收集与分析成本较高；结果性评价便于操作，却容易遮蔽学习过程。如何在二者之间取得平衡，仍需进一步探索。

\section{反思}

回顾两轮行动研究，研究者的主要收获是认识到项目化学习不是一套固定的操作流程，而是一种需要不断调整的设计思维。第一轮实施中，研究者过于关注活动是否完整，忽视了学生的实际困难；第二轮调整后，研究者开始把注意力放在“学生需要什么支架”“哪些证据能够证明学习发生”等问题上。这一转变使项目设计更加贴近学生的学习经验，也使评价更有针对性。

研究者的另一个反思是，地方文化资源的开发不能只依赖教师个人。泉州世界遗产涉及的历史、建筑、民俗与语言知识远超语文教师的专业范围，需要与历史、地理、艺术等学科教师合作，也需要博物馆、文化馆、非遗传承人与社区工作者的支持。未来可以尝试建立跨学科的项目团队，共同开发资源包与教学案例，减轻单一教师的负担。

此外，研究者意识到，项目化学习中的“文化”不应被简化为知识点或展示素材。学生在走访中看到的不仅是古建筑，还有生活在遗产周边的人、正在变化的社区与复杂的保护问题。语文学习应当为学生提供理解这些现象的语言工具与思维方法，而不是把文化遗产包装成一套固定的解说词。只有当学生能够用自己的语言表达真实的观察与思考，文化理解才真正发生。
